% metrics.tex — Metric reference for the explainer evaluation pipeline.
% Compile with: pdflatex -shell-escape metrics.tex  (requires minted + pygments)
\documentclass[11pt,a4paper]{article}

\usepackage[utf8]{inputenc}
\usepackage[T1]{fontenc}
\usepackage{amsmath,amssymb}
\usepackage{booktabs}
\usepackage{hyperref}
\usepackage[most]{tcolorbox}
\usepackage{minted}
\usepackage[margin=2.2cm]{geometry}
\usepackage{enumitem}

\setminted{
  fontsize=\footnotesize,
  breaklines,
  linenos,
  frame=lines,
  framesep=2mm,
}

\newcommand{\code}[1]{\texttt{#1}}
\newcommand{\jsonkey}[1]{\texttt{#1}}
\newcommand{\prob}{\mathbb{P}}
\newcommand{\model}{\Phi}
\newcommand{\Gexp}{G_{\mathrm{exp}}}
\newcommand{\Gcomp}{G \setminus G_{\mathrm{exp}}}

\title{Explainer Metric Reference\\[4pt]
\large Equations, Implementation, and Provenance}
\author{Generated from the \texttt{mprov3\_explainer} codebase}
\date{\today}

\begin{document}
\maketitle
\tableofcontents
\newpage

%% =========================================================================
\section{Notation}
%% =========================================================================

\begin{itemize}[nosep]
  \item $G = (V, E, X^V)$: input graph with $n_V$ nodes, $n_E$ edges, and
        $d$-dimensional node features $X^V \in \mathbb{R}^{n_V \times d}$.
  \item $\model$: trained GNN classifier.
  \item $\model_c(G)$: softmax probability for class~$c$ when $\model$ is
        evaluated on~$G$.
  \item $\hat{y}(G) = \arg\max_c \model_c(G)$: predicted class.
  \item $y$: ground-truth (target) class.
  \item $\Gexp(t)$: hard explanation subgraph at threshold~$t$, induced by the
        top-ranked nodes or edges whose mask value exceeds~$t$.
  \item $\Gcomp(t)$: complement subgraph (everything not in $\Gexp(t)$).
  \item $M \in [0,1]^{n_V}$ or $M \in [0,1]^{n_E}$: soft explanation mask
        after preprocessing (absolute value, min--max normalisation).
  \item $k \in (0,1]$: fraction of mask entries retained by top-$k$
        binarisation (default $k = 0.2$).
  \item $N_t$: number of threshold levels in the Longa sweep (default $100$).
\end{itemize}


%% =========================================================================
\section{Preprocessing pipeline}
\label{sec:preprocessing}
%% =========================================================================

Before any metric is computed, every raw explanation mask passes through three
operations, following the common representation of
Longa et~al.~\cite{longa2025}.

\begin{enumerate}
  \item \textbf{Conversion.}  Edge masks are optionally converted to node
        masks by averaging incident-edge weights per node.
  \item \textbf{Filtering.}  If $\max(M) - \min(M) < \tau$ with
        $\tau = 10^{-3}$, the mask is considered degenerate and the graph is
        marked invalid for headline aggregation.
  \item \textbf{Normalisation.}  Min--max scaling to $[0,1]$:
        \[
          M_{\text{norm}} = \frac{M - \min(M)}{\max(M) - \min(M)}.
        \]
        Degenerate masks (spread $ < \tau$) are replaced by all-zero tensors.
\end{enumerate}

\paragraph{Implementation.}  Local, in
\code{preprocessing.py}:\code{apply\_preprocessing},
\code{normalize\_mask}, \code{edge\_mask\_to\_node\_mask}.

\begin{minted}{python}
# preprocessing.py — normalize_mask
def normalize_mask(mask: torch.Tensor) -> torch.Tensor:
    if mask.numel() == 0:
        return mask
    min_val = mask.min()
    max_val = mask.max()
    if (max_val - min_val).item() < CONSTANT_MASK_TOLERANCE:  # 1e-3
        return torch.zeros_like(mask, device=mask.device)
    return (mask - min_val) / (max_val - min_val)
\end{minted}


%% =========================================================================
\section{Longa metrics (primary — thesis-aligned)}
\label{sec:longa}
%% =========================================================================

These metrics follow \textbf{Longa et~al.}~\cite{longa2025}, Section~5.1.
They operate on \emph{induced subgraphs} (nodes/edges physically removed from
the graph) and return \emph{continuous probability differences}.  They are the
primary metrics for explainer ranking in this codebase.


%% ---------------------------------------------------------------------------
\subsection{Sufficiency (\texorpdfstring{$F_{\mathrm{suf}}$}{Fsuf})}
\label{sec:fsuf}

\paragraph{JSON key.} \jsonkey{mean\_paper\_sufficiency}

\paragraph{Origin.} Longa et~al.~\cite{longa2025}, Eq.\ for $F_{\mathrm{suf}}$
(Section~5.1), based on Debnath et~al.~\cite{debnath2019}.

\paragraph{Equation.}
The raw mask is swept through $N_t - 1$ percentile-based keep fractions
$q_k = 1 - k/N_t$ for $k = 1, \ldots, N_t - 1$.  At each $q_k$, the top-$q_k$
fraction of nodes (or edges, for edge-only explainers) defines the explanation
subgraph $\Gexp(q_k)$.
\begin{equation}
  F_{\mathrm{suf}} = \frac{1}{N_t - 1} \sum_{k=1}^{N_t - 1}
    \bigl[\model_c(G) - \model_c(\Gexp(q_k))\bigr].
  \label{eq:fsuf}
\end{equation}
A \emph{lower} value is better: a sufficient explanation preserves the
target-class probability.  Negative values indicate the subgraph increases
target probability relative to the full graph.

\paragraph{Range.} $[-1, 1]$.  Not clamped before reporting.

\paragraph{Implementation.} Local, in
\code{pipeline.py}:\code{\_paper\_sufficiency\_and\_comprehensiveness}
(node-native path) and
\code{\_paper\_metrics\_from\_edge\_mask} (edge-native path).

\begin{minted}{python}
# pipeline.py — node-native sufficiency (excerpt)
for keep_frac in keep_fractions:  # [0.99, 0.98, ..., 0.01]
    keep_mask = binarize_top_k(node_mask, keep_frac).bool()
    keep_nodes = all_nodes[keep_mask]
    sub_ei, sub_edge_attr, _ = subgraph(
        keep_nodes, edge_index, edge_attr,
        relabel_nodes=True, num_nodes=N, return_edge_mask=True,
    )
    sub_x = x[keep_nodes]
    exp_prob = _predict_target_proba(
        model, x=sub_x, edge_index=sub_ei,
        batch=sub_batch, edge_attr=sub_edge_attr,
        target_class=target_class,
    )
    suf_sum += (full_prob - exp_prob)
return float(suf_sum / len(keep_fractions))
\end{minted}


%% ---------------------------------------------------------------------------
\subsection{Comprehensiveness (\texorpdfstring{$F_{\mathrm{com}}$}{Fcom})}
\label{sec:fcom}

\paragraph{JSON key.} \jsonkey{mean\_paper\_comprehensiveness}

\paragraph{Origin.} Longa et~al.~\cite{longa2025}, Eq.\ for
$F_{\mathrm{com}}$ (Section~5.1), based on Debnath et~al.~\cite{debnath2019}.

\paragraph{Equation.}
Same sweep as $F_{\mathrm{suf}}$, but evaluated on the complement subgraph
$\Gcomp(q_k)$:
\begin{equation}
  F_{\mathrm{com}} = \frac{1}{N_t - 1} \sum_{k=1}^{N_t - 1}
    \bigl[\model_c(G) - \model_c(\Gcomp(q_k))\bigr].
  \label{eq:fcom}
\end{equation}
A \emph{higher} value is better: removing the explanation should damage the
prediction.

\paragraph{Range.} $[-1, 1]$.  Not clamped before reporting.

\paragraph{Implementation.} Local, same functions as
Section~\ref{sec:fsuf}.

\begin{minted}{python}
# pipeline.py — comprehensiveness (excerpt, same loop)
    comp_mask = ~keep_mask
    comp_nodes = all_nodes[comp_mask]
    comp_ei, comp_edge_attr, _ = subgraph(
        comp_nodes, edge_index, edge_attr,
        relabel_nodes=True, num_nodes=N, return_edge_mask=True,
    )
    comp_x = x[comp_nodes]
    comp_prob = _predict_target_proba(
        model, x=comp_x, edge_index=comp_ei,
        batch=comp_batch, edge_attr=comp_edge_attr,
        target_class=target_class,
    )
    com_sum += (full_prob - comp_prob)
\end{minted}


%% ---------------------------------------------------------------------------
\subsection{F1-fidelity (\texorpdfstring{$F_{f1}$}{Ff1})}
\label{sec:ff1}

\paragraph{JSON key.} \jsonkey{mean\_paper\_f1\_fidelity}

\paragraph{Origin.} Longa et~al.~\cite{longa2025}, Section~5.1
(``$f1$-fidelity'').

\paragraph{Equation.}
The harmonic mean of $(1 - F_{\mathrm{suf}})$ and $F_{\mathrm{com}}$, after
clamping both to $[0, 1]$:
\begin{equation}
  F_{f1} = \frac{2\,(1 - F_{\mathrm{suf}}^{\,\mathrm{c}}) \cdot
            F_{\mathrm{com}}^{\,\mathrm{c}}}
           {(1 - F_{\mathrm{suf}}^{\,\mathrm{c}}) +
            F_{\mathrm{com}}^{\,\mathrm{c}}},
  \label{eq:ff1}
\end{equation}
where $F_{\mathrm{suf}}^{\,\mathrm{c}} = \mathrm{clamp}(F_{\mathrm{suf}},\,
0,\,1)$ and $F_{\mathrm{com}}^{\,\mathrm{c}} =
\mathrm{clamp}(F_{\mathrm{com}},\,0,\,1)$.

\paragraph{Range.} $[0, 1]$.  Higher is better.

\paragraph{Implementation.} Local, in
\code{pipeline.py}:\code{\_paper\_f1\_fidelity}.  The clamp is applied only
for the $F_{f1}$ combination; the reported $F_{\mathrm{suf}}$ and
$F_{\mathrm{com}}$ remain raw signed values.

\begin{minted}{python}
# pipeline.py — _paper_f1_fidelity
def _paper_f1_fidelity(Fsuf: float, Fcom: float) -> float:
    if Fsuf != Fsuf or Fcom != Fcom:  # NaN
        return _NAN
    Fsuf_c = max(0.0, min(1.0, Fsuf))
    Fcom_c = max(0.0, min(1.0, Fcom))
    num = 2.0 * (1.0 - Fsuf_c) * Fcom_c
    den = (1.0 - Fsuf_c) + Fcom_c
    return (num / den) if den > 1e-12 else 0.0
\end{minted}


%% =========================================================================
\section{GraphFramEx metrics (diagnostic — collapsed on this dataset)}
\label{sec:graphframex}
%% =========================================================================

These metrics follow \textbf{Amara et~al.}~\cite{amara2022graphframex} and
use PyG's built-in \code{fidelity()} and \code{characterization\_score()}.
They are \emph{class-decision} rates (binary $0$ or $1$ per graph), not
probability differences.

\begin{tcolorbox}[colback=red!5!white, colframe=red!75!black, title=Collapse warning]
On this dataset, these metrics collapse to
$\mathbf{1}[y \neq \text{class}~0]$ for nearly all explainers because:
\begin{enumerate}[nosep]
  \item PyG's \code{fidelity()} evaluates $M \odot X$ on the \emph{full}
        graph topology (not induced subgraphs), so the GNN collapses to its
        prior-dominant class on the heavily degraded input.
  \item The metric is binary per graph ($0$ or $1$), providing insufficient
        resolution to rank explainers.
\end{enumerate}
Changing~$k$ does not fix either structural issue.  Use the Longa metrics
(Section~\ref{sec:longa}) for explainer ranking.
\end{tcolorbox}


%% ---------------------------------------------------------------------------
\subsection{Fidelity\texorpdfstring{$^+$}{+} (top-$k$ binarised)}
\label{sec:fidplus}

\paragraph{JSON key.} \jsonkey{mean\_fidelity\_plus}

\paragraph{Origin.} Amara et~al.~\cite{amara2022graphframex}
(GraphFramEx); implemented via PyG's
\code{torch\_geometric.explain.metric.fidelity}.

\paragraph{Equation (model mode).}
The mask is first binarised at fraction~$k$:
$M_{\mathrm{bin}} = \mathrm{top\text{-}k}(M, k)$.  Then:
\begin{equation}
  \mathrm{Fid}^{+} = 1 - \mathbf{1}\bigl[\hat{y}\bigl(
    (1 - M_{\mathrm{bin}}) \odot X,\; G\bigr) = \hat{y}(G)\bigr],
  \label{eq:fidplus}
\end{equation}
where $\odot$ is element-wise multiplication and the full graph topology
(all edges) is preserved.  Higher is better: removing the explanation
changes the prediction.

\paragraph{Range.} $\{0, 1\}$ per graph; $[0,1]$ as a dataset mean.

\paragraph{Implementation.} \textbf{External} (PyG) +
local binarisation.
\begin{itemize}[nosep]
  \item Binarisation: local, in
        \code{preprocessing.py}:\code{binarize\_top\_k}.
  \item Fidelity call: external, PyG
        \code{torch\_geometric.explain.metric.fidelity}.
  \item Wrapper: local, \code{pipeline.py}:\code{\_compute\_pyg\_fidelity\_top\_k}.
\end{itemize}

\begin{minted}{python}
# pipeline.py — _compute_pyg_fidelity_top_k (excerpt)
def _compute_pyg_fidelity_top_k(explainer, explanation, *, k, ...):
    binarized = _binarize_explanation_top_k(explanation, k=k)
    fid_result = fidelity(explainer, _fidelity_explanation(binarized))
    fid_plus  = float(fid_result[0])
    fid_minus = float(fid_result[1])
    return fid_plus, fid_minus
\end{minted}

\begin{minted}{python}
# preprocessing.py — binarize_top_k (excerpt)
def binarize_top_k(mask: torch.Tensor, k: float) -> torch.Tensor:
    flat = mask.reshape(-1)
    n = flat.numel()
    n_keep = max(1, int(round(k * n + 0.5)))  # ceil
    n_keep = min(n_keep, n)
    threshold = torch.topk(flat, n_keep, largest=True,
                           sorted=False).values.min()
    out = (mask >= threshold).to(mask.dtype)
    return out
\end{minted}


%% ---------------------------------------------------------------------------
\subsection{Fidelity\texorpdfstring{$^-$}{-} (top-$k$ binarised)}
\label{sec:fidminus}

\paragraph{JSON key.} \jsonkey{mean\_fidelity\_minus}

\paragraph{Origin.} Amara et~al.~\cite{amara2022graphframex}.

\paragraph{Equation (model mode).}
\begin{equation}
  \mathrm{Fid}^{-} = 1 - \mathbf{1}\bigl[\hat{y}\bigl(
    M_{\mathrm{bin}} \odot X,\; G\bigr) = \hat{y}(G)\bigr].
  \label{eq:fidminus}
\end{equation}
\emph{Lower} is better: keeping only the explanation should preserve the
prediction.

\paragraph{Range.} $\{0, 1\}$ per graph; $[0,1]$ as a dataset mean.

\paragraph{Implementation.} Same as Section~\ref{sec:fidplus} — both
values are returned by a single call to PyG's \code{fidelity()}.


%% ---------------------------------------------------------------------------
\subsection{Fidelity\texorpdfstring{$^+$}{+} (soft mask)}
\label{sec:fidplussoft}

\paragraph{JSON key.} \jsonkey{mean\_fidelity\_plus\_soft}

\paragraph{Origin.} Amara et~al.~\cite{amara2022graphframex}.

\paragraph{Equation.}
Identical to Eq.~\eqref{eq:fidplus} but with the continuous soft mask $M$
instead of $M_{\mathrm{bin}}$:
\begin{equation}
  \mathrm{Fid}^{+}_{\mathrm{soft}} = 1 - \mathbf{1}\bigl[\hat{y}\bigl(
    (1 - M) \odot X,\; G\bigr) = \hat{y}(G)\bigr].
  \label{eq:fidplussoft}
\end{equation}

\paragraph{Implementation.} \textbf{External} (PyG) — same
\code{fidelity()} call, no binarisation step.  Wrapper:
\code{pipeline.py}:\code{\_compute\_pyg\_fidelity}.


%% ---------------------------------------------------------------------------
\subsection{Fidelity\texorpdfstring{$^-$}{-} (soft mask)}
\label{sec:fidminussoft}

\paragraph{JSON key.} \jsonkey{mean\_fidelity\_minus\_soft}

\paragraph{Origin.} Amara et~al.~\cite{amara2022graphframex}.

\paragraph{Equation.}
\begin{equation}
  \mathrm{Fid}^{-}_{\mathrm{soft}} = 1 - \mathbf{1}\bigl[\hat{y}\bigl(
    M \odot X,\; G\bigr) = \hat{y}(G)\bigr].
  \label{eq:fidminussoft}
\end{equation}

\paragraph{Implementation.} Same as Section~\ref{sec:fidplussoft}.


%% ---------------------------------------------------------------------------
\subsection{Characterisation score (PyG)}
\label{sec:pygchar}

\paragraph{JSON keys.} \jsonkey{mean\_pyg\_characterization},
\jsonkey{mean\_pyg\_characterization\_soft}

\paragraph{Origin.} Amara et~al.~\cite{amara2022graphframex}; implemented
via PyG's \code{characterization\_score}.

\paragraph{Equation.}
Weighted harmonic mean of $\mathrm{Fid}^+$ and $1 - \mathrm{Fid}^-$
with equal weights $w_+ = w_- = 0.5$:
\begin{equation}
  \mathrm{Char} = \frac{1}
    {\dfrac{w_+}{\mathrm{Fid}^+} + \dfrac{w_-}{1 - \mathrm{Fid}^-}}.
  \label{eq:char}
\end{equation}
Higher is better.  Returns $0$ when $\mathrm{Fid}^+ = 0$ or
$\mathrm{Fid}^- = 1$.

\paragraph{Range.} $[0, 1]$.

\paragraph{Implementation.} \textbf{External} (PyG) —
\code{torch\_geometric.explain.metric.characterization\_score}.  Wrapper:
local, \code{pipeline.py}:\code{\_compute\_pyg\_characterization}.

\begin{minted}{python}
# pipeline.py — _compute_pyg_characterization (excerpt)
def _compute_pyg_characterization(fid_plus, fid_minus, *, device, ...):
    return float(characterization_score(
        torch.tensor(fid_plus, device=device),
        torch.tensor(fid_minus, device=device),
        pos_weight=0.5, neg_weight=0.5,
    ).item())
\end{minted}


%% =========================================================================
\section{Diagnostic metrics}
\label{sec:diagnostics}
%% =========================================================================


%% ---------------------------------------------------------------------------
\subsection{Mask spread}
\label{sec:spread}

\paragraph{JSON key.} \jsonkey{mean\_mask\_spread}

\paragraph{Origin.} Project-specific diagnostic, related to the
degenerate-mask filter of Longa et~al.~\cite{longa2025} ($\tau = 10^{-3}$).

\paragraph{Equation.}
\begin{equation}
  \mathrm{spread}(M) = \max(M) - \min(M).
  \label{eq:spread}
\end{equation}

\paragraph{Range.} $[0, \infty)$ in theory; after min--max normalisation,
non-degenerate masks have $\mathrm{spread} = 1$.

\paragraph{Implementation.} Local,
\code{pipeline.py}:\code{\_mask\_spread}.

\begin{minted}{python}
# pipeline.py — _mask_spread
def _mask_spread(mask):
    if mask is None or mask.numel() == 0:
        return 0.0
    return float(mask.max().item() - mask.min().item())
\end{minted}


%% ---------------------------------------------------------------------------
\subsection{Mask entropy (raw)}
\label{sec:entropy}

\paragraph{JSON key.} \jsonkey{mean\_mask\_entropy}

\paragraph{Origin.} Project-specific diagnostic.  Shannon entropy is a
standard information-theoretic measure~\cite{shannon1948}.

\paragraph{Equation.}
The mask is interpreted as an unnormalised probability distribution.
Let $p_i = |M_i| / \sum_j |M_j|$:
\begin{equation}
  H(M) = -\sum_{i:\, p_i > 0} p_i \ln p_i.
  \label{eq:entropy}
\end{equation}
A uniform mask gives $H = \ln N$; a one-hot mask gives $H = 0$.

\paragraph{Range.} $[0,\, \ln N]$ where $N$ is the mask size.  Not
directly comparable across molecules of different sizes.

\paragraph{Implementation.} Local,
\code{pipeline.py}:\code{\_mask\_entropy}.

\begin{minted}{python}
# pipeline.py — _mask_entropy
def _mask_entropy(mask):
    if mask is None or mask.numel() == 0:
        return 0.0
    flat = mask.detach().float().reshape(-1).abs()
    s = float(flat.sum().item())
    if s <= 1e-12:
        return 0.0
    p = flat / s
    nz = p > 0
    return float(-(p[nz] * p[nz].log()).sum().item())
\end{minted}


%% ---------------------------------------------------------------------------
\subsection{Mask entropy (normalised)}
\label{sec:entropy-norm}

\paragraph{JSON key.} \jsonkey{mean\_mask\_entropy\_normalized}

\paragraph{Origin.} Project-specific diagnostic.

\paragraph{Equation.}
\begin{equation}
  H_{\mathrm{norm}}(M) = \frac{H(M)}{\ln |M|},
  \label{eq:entropy-norm}
\end{equation}
where $|M|$ is the number of entries in the mask.  A uniform mask
approaches $1$; a maximally concentrated mask approaches $0$.

\paragraph{Range.} $[0, 1]$.

\paragraph{Implementation.} Local,
\code{pipeline.py}:\code{\_mask\_entropy\_normalized}.

\begin{minted}{python}
# pipeline.py — _mask_entropy_normalized
def _mask_entropy_normalized(mask):
    if mask is None or mask.numel() <= 1:
        return 0.0
    entropy = _mask_entropy(mask)
    max_entropy = math.log(float(mask.numel()))
    if max_entropy <= 0.0:
        return 0.0
    value = entropy / max_entropy
    return max(0.0, min(1.0, float(value)))
\end{minted}


%% ---------------------------------------------------------------------------
\subsection{Degenerate mask count}
\label{sec:degenerate}

\paragraph{JSON key.} \jsonkey{num\_degenerate\_mask}

\paragraph{Origin.} Project-specific diagnostic, operationalising the
filtering step of Longa et~al.~\cite{longa2025}.

\paragraph{Definition.}
\begin{equation}
  n_{\mathrm{degen}} = \bigl|\{i : \mathrm{spread}(M_i) < \tau\}\bigr|,
  \quad \tau = 10^{-3}.
  \label{eq:ndegen}
\end{equation}

\paragraph{Implementation.} Local,
\code{run\_explanations.py}:\code{run\_one\_explainer}.

\begin{minted}{python}
# run_explanations.py
num_degenerate_mask = sum(
    1 for r in results if r.mask_spread < MASK_SPREAD_TOLERANCE
)
\end{minted}


%% =========================================================================
\section{Aggregation}
\label{sec:aggregation}
%% =========================================================================

All per-graph metric values are aggregated into dataset-level means using a
NaN-aware, valid-only protocol:

\paragraph{Valid-only headline means.}
Only graphs with \code{valid=True} (correctly classified, non-degenerate mask,
no NaN in any metric) contribute.  NaN entries within valid graphs are
skipped rather than averaged in as zero.

\paragraph{All-graph diagnostic means.}
The same NaN-skipping mean is computed over \emph{all} explained graphs
(including invalid ones), stored as \jsonkey{mean\_*\_all\_graphs} siblings.

\paragraph{Implementation.} Local,
\code{pipeline.py}:\code{nanmean}, \code{aggregate\_fidelity};
\code{run\_explanations.py}:\code{run\_one\_explainer}.

\begin{minted}{python}
# pipeline.py — nanmean
def nanmean(xs: list[float]) -> float:
    vals = [x for x in xs if x is not None and not math.isnan(x)]
    return (sum(vals) / len(vals)) if vals else float('nan')
\end{minted}


%% =========================================================================
\section{Summary table}
\label{sec:summary}
%% =========================================================================

\begin{table}[ht]
\centering
\footnotesize
\begin{tabular}{@{}llcll@{}}
\toprule
\textbf{Metric} & \textbf{JSON key} & \textbf{Range} &
  \textbf{Impl.} & \textbf{Origin} \\
\midrule
$F_{\mathrm{suf}}$
  & \jsonkey{mean\_paper\_sufficiency}
  & $[-1,1]$
  & Local
  & Longa et~al.~\cite{longa2025} \\
$F_{\mathrm{com}}$
  & \jsonkey{mean\_paper\_comprehensiveness}
  & $[-1,1]$
  & Local
  & Longa et~al.~\cite{longa2025} \\
$F_{f1}$
  & \jsonkey{mean\_paper\_f1\_fidelity}
  & $[0,1]$
  & Local
  & Longa et~al.~\cite{longa2025} \\
\midrule
$\mathrm{Fid}^+$ (top-$k$)
  & \jsonkey{mean\_fidelity\_plus}
  & $\{0,1\}$
  & PyG + local
  & Amara et~al.~\cite{amara2022graphframex} \\
$\mathrm{Fid}^-$ (top-$k$)
  & \jsonkey{mean\_fidelity\_minus}
  & $\{0,1\}$
  & PyG + local
  & Amara et~al.~\cite{amara2022graphframex} \\
$\mathrm{Fid}^+$ (soft)
  & \jsonkey{mean\_fidelity\_plus\_soft}
  & $\{0,1\}$
  & PyG
  & Amara et~al.~\cite{amara2022graphframex} \\
$\mathrm{Fid}^-$ (soft)
  & \jsonkey{mean\_fidelity\_minus\_soft}
  & $\{0,1\}$
  & PyG
  & Amara et~al.~\cite{amara2022graphframex} \\
Char (top-$k$)
  & \jsonkey{mean\_pyg\_characterization}
  & $[0,1]$
  & PyG
  & Amara et~al.~\cite{amara2022graphframex} \\
Char (soft)
  & \jsonkey{mean\_pyg\_characterization\_soft}
  & $[0,1]$
  & PyG
  & Amara et~al.~\cite{amara2022graphframex} \\
\midrule
Spread
  & \jsonkey{mean\_mask\_spread}
  & $[0,\infty)$
  & Local
  & Project-specific \\
Entropy
  & \jsonkey{mean\_mask\_entropy}
  & $[0,\ln N]$
  & Local
  & Shannon~\cite{shannon1948} \\
Entropy (norm)
  & \jsonkey{mean\_mask\_entropy\_normalized}
  & $[0,1]$
  & Local
  & Project-specific \\
$n_{\mathrm{degen}}$
  & \jsonkey{num\_degenerate\_mask}
  & $\mathbb{N}_0$
  & Local
  & Longa et~al.~\cite{longa2025} \\
\bottomrule
\end{tabular}
\caption{Metric summary: JSON key, range, implementation source, and origin.
``Local'' = implemented in \texttt{mprov3\_explainer};
``PyG'' = \texttt{torch\_geometric.explain.metric}.}
\label{tab:summary}
\end{table}


%% =========================================================================
\section{Key differences between Longa and GraphFramEx fidelity}
\label{sec:comparison}
%% =========================================================================

\begin{table}[ht]
\centering
\footnotesize
\begin{tabular}{@{}lll@{}}
\toprule
& \textbf{Longa $F_{\mathrm{suf}}$ / $F_{\mathrm{com}}$}
& \textbf{GraphFramEx $\mathrm{Fid}^+$ / $\mathrm{Fid}^-$} \\
\midrule
Signal type
  & Continuous probability difference
  & Binary class-decision ($0/1$) \\
Perturbation
  & Induced subgraph (remove nodes/edges)
  & Multiplicative masking ($M \odot X$, full topology) \\
Sweep
  & 99 threshold points averaged
  & Single $k$ fraction \\
Sensitivity to explainer
  & High (different masks $\to$ different subgraphs)
  & Low (GNN collapses on degraded input) \\
Origin paper
  & Longa et~al.~\cite{longa2025}
  & Amara et~al.~\cite{amara2022graphframex} \\
\bottomrule
\end{tabular}
\caption{Structural comparison of the two fidelity families.}
\label{tab:comparison}
\end{table}


%% =========================================================================
\begin{thebibliography}{9}

\bibitem{longa2025}
A.~Longa, S.~Azzolin, G.~Santin, G.~Cencetti, P.~Lio, B.~Lepri, and
A.~Passerini,
``Explaining the Explainers in Graph Neural Networks: a Comparative Study,''
\emph{ACM Computing Surveys}, vol.~57, no.~5, Article~120, Jan.~2025.
\textsc{doi}: \href{https://doi.org/10.1145/3696444}{10.1145/3696444}.

\bibitem{amara2022graphframex}
K.~Amara, R.~Ying, Z.~Zhang, Z.~Han, Y.~Shan, U.~Brandes, and S.~Schemm,
``GraphFramEx: Towards Systematic Evaluation of Explainability Methods for
Graph Neural Networks,''
in \emph{Proc.\ Learning on Graphs Conference (LoG)}, 2022.

\bibitem{debnath2019}
J.~DeYoung, S.~Jain, N.~F.~Rajani, E.~Lehman, C.~Xiong, R.~Socher, and
B.~C.~Wallace,
``ERASER: A Benchmark to Evaluate Rationalized NLP Models,''
in \emph{Proc.\ ACL}, 2020.

\bibitem{shannon1948}
C.~E.~Shannon,
``A Mathematical Theory of Communication,''
\emph{Bell System Technical Journal}, vol.~27, pp.~379--423, 1948.

\end{thebibliography}

\end{document}
